\documentclass[10pt,a4paper]{amsart}
\usepackage[margin=19mm]{geometry}
\usepackage{amsmath,amssymb,mathtools}
\usepackage[scr=rsfs,cal=euler]{mathalfa}
\usepackage{xfrac}
\usepackage[unicode,hidelinks,bookmarks=false]{hyperref}
\newcommand{\ie}{i.e.\ }
\newcommand{\cf}{cf.\ }
\newcommand{\resp}{respectively\ }
\newcommand{\Subsection}{\subsection}
\providecommand{\nobreakdash}{\nobreak\hbox{-}}
\setlength{\emergencystretch}{2em}
\multlinegap0pt
\usepackage{bm}
\usepackage{bbm}
\usepackage{dsfont}
\usepackage{xspace}
\usepackage{cancel}
\usepackage{mathtools}
\usepackage{amsthm}
\makeatletter
\@ifundefined{theorem}{\newtheorem{theorem}{Theorem}}{}
\@ifundefined{lemma}{\newtheorem{lemma}[theorem]{Lemma}}{}
\makeatother
\newcommand{\HH}{\mathbf{H}}
\title[Cohomology of the hyperplane complement of a quaternionic re]{Cohomology of the hyperplane complement of a quaternionic reflection group}
\author{Stephen Griffeth, David Guevara}
\date{Source: arXiv:2510.18607v1 (CC BY 4.0)}
\begin{document}
\maketitle

\noindent\emph{Source and licence.} Cohomology of the hyperplane complement of a quaternionic reflection group. Version arXiv:2510.18607v1 (CC BY 4.0), distributed under the Creative Commons Attribution 4.0 International licence, as recorded in the arXiv licence metadata for this exact version and shown on the abstract page https://arxiv.org/abs/2510.18607v1. Full rights notice: licenses/arxiv-2510.18607-CC-BY-4.0.txt.\par\smallskip

\noindent\textit{Editorial scope.} Counted unit: the Lemma whose source label is \texttt{star uniqueness} in the article (source0.tex lines 327--329, source label \texttt{star uniqueness}), delivered together with the article's own complete original proof of it (source0.tex lines 330--346, one \texttt{proof} environment of 17 lines). No mathematics is written, repaired or re-derived: statement and proof are the article's own text line for line. Required context retained uncounted: none. Every definition, display and label of those lines is retained unchanged. Omitted from this package and not redistributed: 1--326, 347--1015. Nothing inside a retained excerpt is omitted or altered: every retained body is delivered exactly as the article's lines are written, so no omission receipt of any kind accompanies this package. Citations: the retained excerpts carry no citation command at all, so no bibliography entry of the article is transcribed and no reference is redistributed. Reference adaptation: none was needed and none was made; every reference inside the delivered unit resolves inside it, so no unresolved reference or citation is produced. Printed slips kept literal: no printed word, symbol, hypothesis or number of the counted statement or of its proof was altered. Layout and class: the article's own class is \texttt{amsart} with packages including fourier, amsfonts, amssymb, graphicx, enumerate, youngtab; the wrapper is the lane's pinned \texttt{amsart} editorial wrapper at 10pt on A4, which restates the theorem-like environments the retained text needs and adds the article's own macro definitions that the retained text needs (\texttt{\textbackslash HH}), with extra wrapper packages (bm, bbm, dsfont, xspace, cancel, mathtools). The delivered numbering is the unit's own. Assets: no retained span contains a figure environment, a graphics-inclusion command, a PDF-page inclusion, a table or a TikZ picture; no image file is copied into this package, the package declares no asset dependency and no image file is redistributed with it. Licence: version arXiv:2510.18607v1 of this article is distributed under the Creative Commons Attribution 4.0 International licence, as recorded in the arXiv licence metadata for this exact version and shown on the abstract page https://arxiv.org/abs/2510.18607v1. A full rights notice is delivered at licenses/arxiv-2510.18607-CC-BY-4.0.txt. Characters: every retained excerpt is pure ASCII, so the delivered engine, XeLaTeX with its default Latin Modern text font, reports no missing character.
\begin{lemma} \label{star uniqueness}
Suppose $k$ and $\ell$ are lines and the angle between them is $\pi/3$. Then there is a unique line $m$ in the plane spanned by $k$ and $\ell$ that is at angle $\pi/3$ to both $k$ and $\ell$.
\end{lemma}

\begin{proof}
Choose basis elements $v_1$ and $v_2$ for $k$ and $\ell$, such that $(v_1,v_1)=2=(v_2,v_2)$. Existence: by scaling $v_1$ we may assume $(v_1,v_2)=-1$. Put $v=-v_1-v_2$. Then 
$$(v,v)=(v_1,v_1)+2(v_1,v_2)+(v_2,v_2)=2-2+2=2 \quad \text{and} \quad (v,v_1)=-(v_1,v_1)-(v_2,v_1)=-2+1=-1=(v,v_2).$$ It follows that the line spanned by $v$ is at angle $\pi/3$ to $k$ and $\ell$. 

Now we prove uniqueness. Suppose $m$ is a line at angle $\pi/3$ to each of $k$ and $\ell$, and choose a basis element $v$ of $m$ such that $(v,v)=2$. By scaling, we may assume 
 $$(v_1,v_2)=-1 \quad \text{and} \quad (v_1,v)=-1.$$ Since the angle between the lines spanned by $v_2$ and $v$ is $\pi/3$ we furthermore have 
 $$|(v_2,v)|=1.$$ Write $v=v_1 a+v_2 b$ for $a,b \in \HH$. We obtain 
$$-1=(v_1,v)=(v_1,v_1 a+ v_2 b)=2a-b \quad \implies \quad b=2a+1$$ and
$$2=(v,v)=(v,v_1)a+(v,v_2)b=-a+(v,v_2)b.$$ Now 
$$(v_2,v)=-a+2b \quad \implies \quad (v,v_2)=-\overline{a}+2\overline{b} \quad \implies \quad 2=-a+(-\overline{a}+2\overline{b})b.$$ Substituting $b=2a+1$ into this gives
$$2=-a+(-\overline{a}+2 (\overline{2a+1}))(2a+1)=-a-\overline{a}-2|a|^2+2(2\overline{a}+1)(2a+1).$$ Simplifying this implies
$$a+\overline{a}=-2|a|^2.$$ 

But also
$$1=|(v_2,v)|^2=|-a+2b|^2=|-a+2(2a+1)|^2=|3a+2|^2.$$ Thus if $a=x+yi+zj+wk$ for real numbers $x,y,z,w$ then $x=\frac{a+\overline{a}}{2}=-|a|^2$ and
$$1=(3x+2)^2+9y^2+9z^2+9w^2=9(x^2+y^2+z^2+w^2)+12x+4=9|a|^2-12|a|^2+4.$$ It follows that $|a|=1$, and hence $x=-1$ and finally $a=-1$, $b=2a+1=-1$. 
\end{proof}

\end{document}
